\section{Conclusiones y recomendaciones}
    \begin{itemize}
        \item El protocolo de contraste confirmó que el sector determina el consumo
        energético de forma inequívoca, ya que el ANOVA atribuye el 83,9\,\% de la
        variabilidad a ese factor ($F = 775{,}99$; $p < 0{,}001$) y las tres
        comparaciones de Tukey resultan significativas, con diferencias de 244,3,
        456,4 y 700,7 kWh entre pares de sectores.

        \vspace{10pt}

        \item La verificación de supuestos no fue un trámite sino una decisión de
        método, puesto que el rechazo de la homogeneidad de varianzas obligó a
        emplear la corrección de Welch en la comparación de medias; omitir ese paso
        habría producido un valor $p$ calculado bajo un supuesto que los datos
        contradicen.

        \vspace{10pt}

        \item El incumplimiento de la homocedasticidad se resolvió contrastándolo y
        no advirtiéndolo, ya que repetir el análisis con el ANOVA de Welch
        ($F = 830{,}22$) y el post-hoc de Games-Howell reprodujo íntegramente la
        decisión del procedimiento clásico; la elección del método altera la
        precisión de cada intervalo, estrechándolo entre los sectores de menor
        varianza y ensanchándolo en los que involucran al Industrial, pero no la
        conclusión.

        \vspace{10pt}

        \item Reportar la magnitud junto al valor $p$ resultó indispensable para
        interpretar dos resultados que el estadístico presenta de forma engañosamente
        similar, dado que la diferencia entre sectores corresponde a un efecto
        enorme ($d = 3{,}33$; $\omega^2 = 0{,}838$) mientras que la correlación
        global no alcanza siquiera el umbral de $r_{\min} = 0{,}161$ detectable con
        potencia del 80\,\%, lo que descarta que su no significancia se deba a un
        tamaño muestral insuficiente.

        \vspace{10pt}

        \item El hallazgo central es que una prueba de hipótesis puede no rechazar
        $H_0$ aunque el efecto exista, ya que la correlación temperatura-consumo
        resulta no significativa a nivel agregado ($r = 0{,}063$; $p = 0{,}277$)
        pese a haberse introducido por diseño, y sí lo es dentro del sector
        Residencial ($r = 0{,}595$; $p < 0{,}001$).

        \vspace{10pt}

        \item El mecanismo de esa discrepancia quedó cuantificado mediante la
        relación $r = b_1 s_T / s_Y$, dado que la pendiente sobrevive a la
        agregación ($b_1 = 3{,}42$ frente al valor de diseño de 4) mientras la
        desviación estándar del consumo se triplica al mezclar sectores, de 40,6 a
        317,5 kWh, diluyendo la correlación sin alterar el efecto subyacente.

        \vspace{10pt}

        \item La visualización fue el instrumento que hizo visible esa estructura,
        puesto que la nube agregada de la Figura \ref{fig_regresion} muestra tres
        franjas horizontales atravesadas por una recta que no describe a ninguna,
        mientras que la figura interactiva de la Figura \ref{fig_plotly} revela tres
        pendientes positivas paralelas que ningún estadístico agregado permite
        anticipar.

        \vspace{10pt}

        \item La equivalencia entre Python y R, con coincidencia dígito a dígito en
        Shapiro-Wilk, la t de Welch, el ANOVA, Tukey, la regresión, los tamaños de
        efecto y ambos contrastes robustos, confirma que los resultados dependen de
        los datos y no de la herramienta; el acuerdo es especialmente exigente en
        Games-Howell y en la potencia, programados de cero en cada entorno a partir
        de sus fórmulas y no heredados de una biblioteca común.

        \vspace{10pt}

        \item Las nueve figuras se construyeron dos veces, con Matplotlib, seaborn y
        Plotly en Python y con ggplot2 en R, y cada una lleva impreso el estadístico
        que representa; las dos versiones interactivas y el tablero de cuatro paneles
        permiten además condicionar por sector sin recalcular nada, que es la
        operación exacta que separa la conclusión correcta de la equivocada.
    \end{itemize}

    \vspace{10pt}

    A partir de lo anterior se derivan tres recomendaciones para la toma de
    decisiones:

    \vspace{10pt}

    \begin{itemize}
        \item \textbf{Segmentar antes de promediar.} La media global de 501,98 kWh
        no describe a ningún cliente real, ya que se sitúa entre el sector Comercial
        y el Residencial sin representar a ninguno. Toda tarifa, meta de eficiencia
        o proyección de demanda debe construirse por sector, decisión que los
        intervalos de Tukey respaldan al excluir el cero con márgenes amplios y que
        descarta cualquier esquema tarifario de dos escalones.

        \vspace{10pt}

        \item \textbf{Incorporar la temperatura como variable de planeación
        únicamente en el segmento residencial.} Es el único donde la relación es a
        la vez significativa y de magnitud apreciable ($r = 0{,}595$), lo que la
        vuelve útil para anticipar picos de demanda en las regiones Caribe y
        Pacífica; extender esa conclusión al conjunto agregado carecería de
        sustento.

        \vspace{10pt}

        \item \textbf{Acompañar todo contraste con su representación gráfica antes
        de aceptar un resultado no significativo.} Este trabajo muestra que un valor
        $p$ elevado admite dos lecturas incompatibles, la ausencia de efecto y el
        enmascaramiento por agregación, que el estadístico no distingue y la figura
        sí. En conjuntos reales, donde la variable responsable de la agregación
        puede no estar registrada, esa inspección visual es la única salvaguarda
        disponible frente a una conclusión equivocada.
    \end{itemize}
